\section{Executed experiments and their interpretation}

\subsection{Environment and reproducibility}

The recorded run used Python 3.13.5 on Linux, NumPy 2.3.5, and SciPy 1.17.0.
The sparse polynomial representation and all acceptance checks use Python's
standard library, chiefly exact \code{Fraction} arithmetic. SciPy is used only
for LP proposal generation. SymPy was available but is not needed by the
supplied prototypes.

The recorded seeds are 420 for recurrence cases, 421 for witnesses, 422 for
square certificates, and 423 for Bernstein cases. The test suite uses separate
fixed seeds for its randomized regression checks. The archive includes every
accepted certificate and its stated problem, not just aggregate counts.

No Lean executable was available, and an attempted toolchain acquisition did
not succeed. Consequently no Lean tactic was run, no imported Lean library
combination was compiled, and no kernel-checking time was measured. The file
\code{lean/ReplayExamples.lean} is explicitly labeled uncompiled and is excluded
from all success counts.

\subsection{Results}

\begin{center}
\begin{tabular}{@{}lrrr@{}}
\toprule
Experiment family & Cases & Exact replays & Time (s) \\
\midrule
List induction and generalization & 6 & 6 & 0.005 \\
Polynomial additive recurrences & 65 & 65 & 0.241 \\
Affine witnesses with Farkas certificates & 63 & 63 & 0.352 \\
Finite-dictionary square certificates & 40 & 40 & 0.648 \\
Bernstein box certificates & 41 & 41 & 0.067 \\
\midrule
Total certificate cases & 215 & 215 & \\
\bottomrule
\end{tabular}
\end{center}

The times are one local Python run, rounded from
\code{results/results.json}. Each family includes its search and immediate
checks; these are not independent hardware-normalized benchmark timings or
predictions of Lean performance. The separate replay run verifies all 215
stored certificates without redoing search.

The list family contains four supplied seed statements with searched proofs,
the synthesized accumulator specification, and the original specialized
specification. Its six proof traces contain respectively 6, 3, 8, 10, 9, and
5 rewrites. These are small proofs, as expected for a deliberately restricted
language.

The recurrence family comprises five named power sums and 60 seeded random
polynomial increments with rational coefficients and rational initial values.
It required 312 candidates in total. The exact step checks generated 247
counterexamples before the 65 final certificates were accepted.

The witness family comprises three named inequality problems, 40 problems
whose equalities force affine witnesses, and 20 constructed polyhedral
inequality problems. These exercise both the coefficient synthesis and the
nonnegative multiplier checks. They are not an evaluation on arbitrary
quantified formulas.

The 40 square-certificate problems are \emph{constructed from the same finite
dictionary used by the search}. Their 100\% success rate therefore tests
implementation and rational reconstruction on that designed distribution;
it must not be presented as general nonlinear proving coverage.

The 41 Bernstein problems are strictly positive low-degree polynomials made
from squares plus positive rational constants. Root-only checking certifies
25. Subdivision certifies all 41, using 82 leaves in total and a maximum split
depth of four. A full SOS solver or another tactic may also solve these
problems; no comparative claim is made about them.

\Needspace{10\baselineskip}
\subsection{Ablations that demonstrate the mechanisms}

\begin{longtable}{@{}>{\raggedright\arraybackslash}p{0.32\textwidth}>{\raggedright\arraybackslash}p{0.29\textwidth}>{\raggedright\arraybackslash}p{0.33\textwidth}@{}}
\toprule
Problem / mechanism & Restricted attempt & Enabled mechanism \\
\midrule
\endhead
Accumulator reverse & Induction with empty accumulator fixed: no proof. & Changing-argument generalization and typed RHS enumeration: checked proof. \\
Universal strip $x+1\le y\le x+2$ & Constant witness template: no certificate. & Affine template: $y=x+1$, exactly checked. \\
41 positive box polynomials & Root Bernstein coefficients: 25 certified. & Rational subdivision: 41 certified. \\
Triangular-number interpolation trap & First four values agree with a false formula. & Exact induction-step identity rejects it. \\
Diagonal square $(x-y)^2$ & Bernstein search at depth four: unknown. & Square certificate is available directly. \\
\bottomrule
\end{longtable}

These ablations answer the question ``does the proposed mechanism do something
that its restricted precursor does not?'' They do not answer ``how many more
Mathlib goals does \Forge{} solve than \Grind{}?'' That second question requires
an implemented Lean tactic and a controlled comparison.

\subsection{Negative and adversarial tests}

All 31 unittest methods passed. Several contain multiple seeded trials,
including exact polynomial arithmetic, Bernstein reconstruction, recurrence
identities, and semantic cross-checks of the list equations. The suite also
rejects a false reverse identity, wrong-tail induction matching, an unknown
or self-referential rule, a corrupt rewrite path, and ill-sorted terms.

Arithmetic checks reject negative square weights, certificates attached to a
changed target, omitted domain constraints, an incorrect witness coefficient,
inexact witness-certificate values, and negative slacks. Partition checks
reject missing children, boundary split points, and invalid axes. Strict
positivity and weak nonnegativity are tested separately.

The Motzkin polynomial
\[
x^4y^2+x^2y^4+1-3x^2y^2
\]
is a deliberate unsuccessful plain-square-dictionary case. No certificate is
reported. The test is a regression boundary for this solver, not evidence
that the polynomial is negative. Similarly, insufficient recurrence degree
or insufficient grammar size returns no proof, and a depth-limited Bernstein
failure remains unknown.

\subsection{What the evidence does and does not establish}

It establishes that the included programs ran, that the reported certificates
passed their specified exact Python checks, and that the listed negative
cases were rejected. Together with the mathematical arguments, it supports
the feasibility of several workers and the usefulness of generalization,
certificate synthesis, and subdivision in their stated fragments.

It does not establish correctness of a Lean implementation that has not been
written, completeness outside those fragments, robustness of a dependent-type
reifier, improved performance on a representative corpus, or a measured
advantage over current \Grind{}. There is also shared substrate between
search and replay: both use the same term or polynomial representation.
``Search-free replay'' is not the same as an independently implemented formal
verification of that representation.

\Needspace{9\baselineskip}
\section{How to demonstrate genuine superiority over \Grind{}}

\subsection{Freeze the comparison, including its theorem knowledge}

Use a pinned Lean release, a compatible pinned Mathlib revision, fixed imported
modules, and recorded backend versions. Preserve each source goal and its
original hypotheses. Avoid accidentally making the tested theorem itself,
a renamed duplicate, or a future theorem dependent on it available through
library retrieval. This is especially important when benchmarks are extracted
from an already compiled mathematical library.

Compare ordinary \Grind{}, \Grind{} with explicit relevant equations/lemmas,
and the current shared-state bit-vector route where applicable. Also compare
Aesop, a simple tactic portfolio, and appropriate specialized tools such as
SOS, Duper, or Lean-SMT. A system that wins only against an artificially weak
baseline has not demonstrated the desired improvement.

For structural problems, distinguish three input conditions: definitions only;
a common relevant helper-lemma set; and an oracle-assisted set containing the
handwritten proof's helper lemmas. The gap between the first two conditions
measures lemma acquisition, while the last helps separate discovery failures
from local proof failures.

\subsection{The benchmark families}

The most informative families are recursive data-structure theorems needing
induction or changing accumulators; polynomial cost and invariant proofs;
quantified statements needing new arithmetic terms; constructive existential
specifications; mixed guarded arithmetic and rewriting; and higher-order or
indexed-datatype obligations.

Add a large control set of goals already handled well by \Grind{}, including
its arithmetic and bit-vector strengths. The control set detects whether the
new controller wastes resources on familiar cases. DTTBench, introduced with
Canonical-min, is relevant additional material for the typed-inhabitation
worker, but does not replace a Lean integration benchmark.\cite{canonicalmin}

Hold out whole problem families when tuning heuristic weights. Renaming
variables or slightly changing constants in a training problem does not
create an independent test set for lemma-discovery heuristics.

\subsection{Metrics and ablations}

Report solved fraction, wall time, Lean heartbeats where meaningful, peak
memory, proof/certificate size, reconstruction time, and kernel checking time.
Record failed translations, unsupported fragments, timeouts, rejected
certificates, and residual reconstruction goals separately. A solver verdict
with an unclosed proof hole is not a solved benchmark.

Run at least the following ablations: no induction; induction without
generalization; no helper-lemma discovery; no constrained instantiation;
constant-only witnesses; no nonlinear certificate worker; and stateless local
solver restarts versus the incremental adapter. Include the simple portfolio
with the same total budget to measure the controller's interactions rather
than merely its larger collection of tools.

The principal evidence for ``significantly more capable'' should be a
substantial held-out solved-set increase in structural and synthesis-heavy
families, while retaining most control-set successes and producing acceptable
proof-checking costs. A speed claim requires repeated controlled runs; a
capability claim requires concrete additional checked goals. Neither follows
from the present prototype counts alone.

\section{Assessment and recommended starting point}

The most promising first implementation is \emph{not} a complete new SMT
solver and not a rewrite of \Grind{}. It is a narrow Lean controller that
combines automatic generalization and induction with an exact witness and
recurrence-certificate pipeline. This already changes the shapes of proofs
that can be found automatically. It also provides clear interfaces for adding
stronger nonlinear and quantified reasoning afterward.

The highest-risk engineering tasks are preserving dependent contexts through
generalization, maintaining valid state across local solver restarts or
snapshots, producing compact proof reconstruction, and keeping candidate
lemmas out of the proved state. The numerical search itself is comparatively
easy to sandbox because an exact certificate can reject a bad proposal.

The Python experiments give a concrete foothold: a synthesized accumulator
lemma, exact recurrence invariants, parametric witnesses with proof
multipliers, and complementary polynomial certificates. The proposed Lean
architecture explains how these become cooperating proof workers rather
than isolated demos. That is a technically plausible route to a tactic
substantially more useful than unassisted local saturation on the targeted
goal families. Its eventual superiority should be earned by the pinned,
proof-checked benchmark program just described.

\appendix
\section{Artifact inventory and commands}

The archive is organized as follows:
\begin{lstlisting}
article/
  forge.tex, 01-architecture.tex, 02-algorithms.tex
  03-integration.tex, 04-evaluation.tex, references.tex
  forge.pdf
prototype/
  forge_proto/{poly,linear,induction,recurrence,
               witness,positivity}.py
  tests/test_forge.py
  run_experiments.py
  replay_certificates.py
  requirements.txt, README.md
results/
  results.json, experiments.log, tests.log, replay.log
  certificates/{induction,recurrences,witnesses,
                sos,bernstein}.json
lean/
  ReplayExamples.lean
README.md, SOURCES.md
\end{lstlisting}

From \code{prototype/}, reproduce the executed work with:
\begin{lstlisting}
python -m pip install -r requirements.txt
python -m unittest discover -s tests -v
python run_experiments.py --output ../results
python -S replay_certificates.py ../results/certificates
\end{lstlisting}
The \code{-S} option disables third-party site-package loading. The last
command succeeds because checking does not require the numerical proposer.
It reports 6 induction, 65 recurrence, 63 witness, 40 square, and 41 Bernstein
certificates, totaling 215.

Build this article from \code{article/} with two runs of LuaLaTeX:
\begin{lstlisting}
lualatex -interaction=nonstopmode -halt-on-error forge.tex
lualatex -interaction=nonstopmode -halt-on-error forge.tex
\end{lstlisting}
The document uses Linux Libertine O, Lato, and DejaVu Sans Mono installed as
system fonts. Font files are not included. Another available font family can
be selected by changing the three \code{fontspec} declarations.

\section{Certificate representations}

\subsection{Polynomials}

A polynomial stores a positive variable count and a sorted list of nonzero
rational coefficients paired with nonnegative exponent tuples. Rational
coefficients are serialized as exact strings such as \code{"3/7"}. Canonical
construction rejects malformed exponent vectors, zero entries, and duplicate
or unsorted serialized terms.

The variable table binding these indices to Lean expressions is a separate
production concern. The Python arithmetic certificates use fixed mathematical
variables and do not implement that Lean source map.

\subsection{Induction traces}

An induction certificate contains an equation, its induction variable, and
two equality traces. Each equality trace has a sequence of rewrites for its
left and right sides. A rewrite is a pair \code{(rule_name, subterm_path)}.
The checker derives the base/step contexts and the only permitted induction
hypothesis; these are not arbitrary claims supplied by the certificate.

The fixture stores certificates in dependency order. Replay installs each
one only after checking it against defining equations and earlier installed
lemmas. The experiment's seed statements therefore do not enter as axioms.

\subsection{Affine and square certificates}

An affine witness certificate stores one affine expression per output,
one nonnegative multiplier row per target inequality, and one nonnegative
slack per target inequality. Verification is exact coefficient equality in
\eqref{eq:farkas-linear}--\eqref{eq:farkas-constant}.

A square certificate stores weighted square roots with an optional index of
a source nonnegative constraint. It is checked against the separately supplied
target and constraint list. The current format has no arbitrary equality
multipliers; adding those is a straightforward extension of the identity
checker, but not an implemented feature claimed by the experiments.

\subsection{Bernstein trees}

A leaf contains only its kind. A split contains an axis, exact split point,
and two subtrees. The checker derives all bounds from the root domain and
recomputes all leaf coefficients. Node/depth limits are operational rejection
limits. They are not mathematical assumptions about positivity.

\subsection{A production interchange format}

A production format should additionally carry a schema version, source-goal
identifier, environment/encoding revision, variable-map identifier, required
assumptions, and claimed result kind. The Lean process resolves these
identifiers against its own state rather than accepting oracle-provided
expressions as authoritative. Resource limits should bound integer bit sizes,
monomial counts, degrees, tree nodes, and proof depth before expensive checking.

\section{Known limitations of the delivered code}

The prototypes do not implement the full obligation scheduler, actual Lean
snapshots, native \Grind{} state access, dependent typing, general quantifier
instantiation, higher-order synthesis, or external SMT proof reconstruction.
The square search is a finite linear dictionary, not semidefinite programming.
Affine synthesis is over rational/real non-strict polyhedral constraints,
not mixed-integer or arbitrary quantified arithmetic. The recurrence solver
handles additive polynomial increments only.

The exact checkers are ordinary Python programs. They have executable
regression tests and elementary soundness arguments for their formats, but
are not themselves formally verified. Several share representation and
matching routines with their search counterparts. This is appropriate for a
prototype and is precisely why the Lean implementation must reconstruct
proofs or use a checker with a proved soundness theorem.

The Lean examples are uncompiled illustrations with no \code{sorry}
placeholders. Absence of \code{sorry} is not evidence that they compile, and
no such claim is made. The package is an actionable research/design artifact
with tested algorithmic components, not a released drop-in replacement for
\Grind{}.
